% Header: General study / work / project template
%%%%%%%%%%%%%%%%%%%%%%%%%%%%%%%%%%%%%%%%%%%%%%%%%%%%%%%%%%%%%%%%%%%

\documentclass[11pt,a4paper]{scrartcl}

% Layout
\usepackage[margin=2.5cm]{geometry}
\usepackage[onehalfspacing]{setspace}

% General packages
\usepackage{graphicx}
\usepackage{enumitem}
\usepackage{tcolorbox}
\tcbuselibrary{breakable}
\usepackage[breaklinks=true,colorlinks=true,linkcolor=blue,urlcolor=blue,citecolor=blue]{hyperref}
\usepackage{amsmath,amsthm,amssymb,mathtools}
\usepackage{braket}
\usepackage{icomma}
\usepackage[T1]{fontenc}
\usepackage[utf8]{inputenc}

% Language: choose one
% \usepackage[ngerman]{babel}
\usepackage[english]{babel}

% Units / tables / floats / references
\usepackage[locale = UK,space-before-unit=true,per-mode = symbol]{siunitx}
\usepackage{booktabs,multirow}
\usepackage{placeins}
\usepackage[natbib,abbreviate=true,doi=false,style=numeric-comp,giveninits=true,sorting=none]{biblatex}
\usepackage{csquotes}

% Optional bibliography
% \addbibresource{references.bib}

% Graphics paths
\graphicspath{{Bilder/} {images/} {figures/} {chapters/}}

% Optional math shortcuts
\newcommand{\R}{\mathbb{R}}
\newcommand{\N}{\mathbb{N}}
\newcommand{\Q}{\mathbb{Q}}
%\newcommand{\set}[1]{\left\{ #1 \right\}}
\newcommand{\abs}[1]{\left| #1 \right|}
\newcommand{\norm}[1]{\left\lVert #1 \right\rVert}
\newcommand{\ol}[1]{\overline{#1}}

% Optional units
\DeclareSIUnit{\dBm}{dBm}

% Box styles
\newtcolorbox{calculationbox}[1][]{
    title=Calculation,
    breakable,
    #1
}

\newtcolorbox{solutionbox}[1][]{
    title=Solution,
    breakable,
    #1
}

\newtcolorbox{answerbox}[1][]{
    title=Answer,
    breakable,
    #1
}

\newtcolorbox{notebox}[1][]{
    title=Notes,
    breakable,
    #1
}

\newtcolorbox{sidecalcbox}[1][]{
    title=Side Calculation,
    breakable,
    #1
}

%%%%%%%%%%%%%%%%%%%%%%%%%%%%%%%%%%%%%%%%%%%%%%%%%%%%%%%%%%%%%%%%%%%
\begin{document}

    \title{Theoretical Physics 2 - Blatt 06}
    \author{Tobias Laser}
    \date{\today}
    \maketitle
    \vfill
    \thispagestyle{empty}

    \pagenumbering{arabic}

    % ================================================================
    \paragraph{19. Zeitentwicklung und Messungen eines Spins mit $s = 1 / 2$}:
    
    Wir betrachten ein Spin-$\frac{1}{2}$-Teilchen, das sich zum Zeitpunkt $t = 0$ im  Zustand $\ket{\psi(t = 0)} = \ket{\downarrow}$ befindet. Die Zeitentwicklung wird durch den Hamiltonoperator
    \begin{align*}
        \operatorname{H} = \frac{\hbar \omega}{2} \sigma_x,
    \end{align*} 
    beschrieben.
    \begin{enumerate}
        \item Es wird zum Zeitpunkt $t_1 = \frac{\pi}{2 \omega}$ eine Messung des Operators $\sigma_z$ durchgeführt. In welchem Zustand befindet sich das Teilchen zu diesem Zeitpunkt? Welche Messergebnisse sind möglich, und mit welchen Wahrscheinlichkeiten tretten sie auf?
        \item Direkt nach der Messung entwickelt sich das System für ein Zeitintervall $\Delta t = \frac{\pi}{2\omega}$ weiter. Zum Zeitpunkt $t_2 = t_1 + \delta t = \frac{\pi}{\omega}$ wird erneut $\sigma_z$ gemessen. Abhängig vom Messergebnis aus Teilaufgabe (1): In welchem Zustand befindet sich das Teilchen zum Zeitpunkt $t_2$, und mit welchen Wahrscheinlichkeiten treten die jeweiligen Messergebnisse auf?
        \item Im Gegensatz zu den vorherigen Teilaufgaben entwickelt sich das Teilchen nun ungestört (d.h. ohne Messung bei $t=t_1$) bis zum Zeitpunkt $t_2$. In welchem Zustand befindet sich das Teilchen zu diesem Zeitpunkt? Welche Wahrscheinlichkeiten ergeben sich bei einer Messung von $\sigma_z$ zum Zeitpunkt $t_2$? Begründen Sie die Unterschiede zu Teilaufgabe (2).
    \end{enumerate}

    \paragraph{20. System mit Drehimpuls $l=1$}:

    Betrachten Sie ein System mit Drehiumpuls $l=1$. Die drei Vektoren $\ket{+1}, \ket{0}, \ket{-1}$, Eigenvektoren von $\operatorname{L}_z$ zu den Eigenwerten $\hbar, 0, -\hbar$, bilden eine Basis des zugehörigen Zustandsraumes. Für diese Eigenvektoren gilt
    \begin{align*}
        \operatorname{L}_\pm\ket{m} = \hbar \sqrt{2} \ket{m \pm 1}\\
        \operatorname{L}_+\ket{1} = \operatorname{L}_- \ket{- 1} = 0\\
    \end{align*}
    Dijeses System besitzt ein elektrisches Quadrupolmoment und befinde sich in einem elektrischen Feldgradienten. Der Hamiltonoperator ist gegeben durch
    \begin{align*}
        \operatorname{H} = \frac{\omega_0}{\hbar} (\operatorname{L}_u^2 - L_v^2),
    \end{align*}
    wobei $\operatorname{L}_u$ und $\operatorname{L}_v$ die Komponenten von $\operatorname{L}$ entlang der Diagonalen $u$ und $v$ der $xz$-Ebende sind; $\omega_0$ ist eine reelle Konstante.
    \begin{enumerate}
        \item Geben Sie die Matrixdarstellung für $\operatorname{H}$ in der Basis $\{\ket{+1}, \ket{0}\ket{-1}\}$ an. Wie lauten die stationären Zustände des Systems und ihre Energieeigenwerte? Schreiben Sie diese Zustände als $\ket{E_1}, \ket{E_2}, \ket{E_3}$, mit $E_1 \geq E_2 \geq E_3$.
        \begin{calculationbox}
            \begin{calculationbox}
                \begin{align*}
                    \operatorname{H}&=\frac{\omega_0}{\hbar} \left(
                        \operatorname{}
                    \right)
                    \\
                    &=
                    \frac{\omega_0}{\hbar} \left( 
                        \hbar
                        \begin{pmatrix}
                            1&0&0\\
                            0&0&0\\
                            0&0&-1\\
                        \end{pmatrix}
                        \frac{\hbar}{\sqrt{2}}
                        \begin{pmatrix}
                            0&1&0\\
                            1&0&1\\
                            0&1&0\\
                        \end{pmatrix}
                        + \frac{\hbar}{\sqrt{2}}
                        \begin{pmatrix}
                            0&1&0\\
                            1&0&1\\
                            0&1&0\\
                        \end{pmatrix}
                        \hbar
                        \begin{pmatrix}
                            1&0&0\\
                            0&0&0\\
                            0&0&-1\\
                        \end{pmatrix}
                    \right)\\
                    &=
                    \frac{\hbar \omega_0}{\sqrt{2}}
                    \left(
                        \begin{pmatrix}
                            0&1&0\\
                            0&0&0\\
                            0&-1&0\\
                        \end{pmatrix}
                        +
                        \begin{pmatrix}
                            0&0&0\\
                            1&0&-1\\
                            0&0&0\\
                        \end{pmatrix}
                    \right)\\
                    &=
                    \frac{\hbar \omega_0}{\sqrt{2}}
                    \begin{pmatrix}
                        0&1&0\\
                        1&0&-1\\
                        0&-1&0\\
                    \end{pmatrix}
                \end{align*}
            \end{calculationbox}
        \end{calculationbox}
        \item Zur Zeit $t=0$ sei das System im Zustand
        \begin{align*}
            \ket{\psi(0)} = \frac{1}{\sqrt{2}}[\ket{+1} - \ket{-1}].
        \end{align*}
        Wie lautet der Zustandsvektor $\ket{\psi(t)}$ Zur Zeit $t$ werde $\operatorname{L}_z$ gemessen; wie groß sind die Wahrscheinlichkeiten der verschiedenen möglichen Ergebnisse?
        \item Berechen Sie die Erwartungswerte $\braket{\operatorname{L}_x(t)}, \braket{\operatorname{L}_y(t)}$ und $\braket{\operatorname{L}_z(t)}$. Welche Bewegung führt der Vektor $\braket{\operatorname{L}(t)}$ aus? Skizzieren Sie die Trajektorie von $\braket{\operatorname{L}(t)}$.
        \item Zur Zeit $t$ werde eine Messung von $\operatorname{L}_z^2$ durchgeführt.
        \begin{enumerate}
            \item Gibt es Zeiten, wo nur ein einziges Ergebnis möglich ist?
            \item Nehmen Sie an, dass diese MEssung das Ergebnis $\hbar^2$ ergeben hat. In welchem Zustand befindet sich das System unmittelbar nach der Messung? Beschreiben Sie ohne Rechnung die nachfolgende entwicklung.
        \end{enumerate}
    \end{enumerate}
\end{document}